\documentclass[12pt]{article}
\usepackage[utf8]{inputenc}
\usepackage[T1]{fontenc}
\usepackage{amsmath,amssymb}
\usepackage{geometry}
\geometry{margin=2.5cm}

\begin{document}

\noindent $xR_H^L y \Leftrightarrow (x,y)\in R_H^L$.

\noindent $x\equiv_L y \pmod H \Leftrightarrow x^{-1}y\in H$.

\bigskip
\noindent $x\equiv_L y \pmod H \Leftrightarrow y\equiv_L x \pmod H$ $(\forall)\, x,y\in G$.

\noindent $x\equiv_L y \pmod H$, $y\equiv_L z \pmod H \Rightarrow x\equiv_L z \pmod H$

\bigskip
\noindent $x\equiv_L y \pmod H \Rightarrow zx\equiv_L zy \pmod H$. \ $(\forall)\, x,y\in G$.\footnote{``z'' is unclear in the original (it could look like the digit ``2''), but the meaning matches the left-compatibility established on the previous page ($xRy\Rightarrow zxRzy$).}

\bigskip
\noindent\underline{Remark}: If $R$ is an equivalence relation on $G$ compatible with left multiplication $\Rightarrow$ $\exists\, H\leq G$ such that $R=R_H^L$

\[
H \overset{\text{def}}{=} [e]_R = \{x\in G \mid xRe\}.
\]

\noindent $x,y\in H \overset{?}{\Rightarrow} xy\in H$.

\noindent $xRe$

\noindent $yRe$

\noindent $\Rightarrow xyRe \Rightarrow xy\in H$.

\bigskip
\noindent We show that $R=R_H^L$.

\noindent Let $(x,y)\in R \Leftrightarrow xRy \Leftrightarrow eRx^{-1}y \Leftrightarrow x^{-1}yRe \Leftrightarrow x^{-1}y\in H$

\noindent $\Leftrightarrow (x,y)\in R_H^L$.

\bigskip
\noindent\underline{Prop}: $G$ a group and $H\leq G$; and $a\in G$ such that

\[
[a]_{R_H^L} = aH \overset{\text{def}}{=} \{ah \mid h\in H\}
\]

\noindent = the left coset with representative $a$ modulo the subgroup $H$.

\bigskip
\noindent Suppose $x\in[a]_{R_H^L} \Rightarrow x\equiv_L a \pmod H \Rightarrow a\equiv_L x \pmod H$

\noindent $\Leftrightarrow a^{-1}x\in H \Leftrightarrow x=ah,\ h\in H$.

\bigskip
\noindent Let $\{a_i\}_{i\in I}$ be a right transversal of $H$ in $G$.

\[
\{a_i\}_{i\in I} \xrightarrow{\ \xi\ } \{a_i^{-1}\}_{i\in I} \qquad
\begin{aligned}
&(\forall)\, i\neq j \Rightarrow a_i \not\equiv_L a_j \pmod H\\
&(\forall)\, x\in G,\ \exists!\, i\in I \text{ such that } x\equiv_L a_i \pmod H.
\end{aligned}
\]

\noindent $\xi$ is a bijection.

\[
\text{Card}(I) = [G:H]
\]

\noindent = the index of $H$ in $G$. (the number of left (right) residue classes modulo the subgroup $H$).

\end{document}
